\begin{table}[ht]
\centering
\begin{tabular}{lccccccccc}
\toprule
 & Persistent:
 constant & Persistent:
 decreasing strength & Persistent:
 increasing strength & Lost:
 rewiring & Gained:
 rewiring & Lost:
 one transient sp. & Lost:
 both transient sp. & Gained:
 one transient species & Gained:
 both transient sp. \\
\midrule
Persistent:
 constant & -- & 1.000 & 1.000 & 1.000 & 1.000 & 1.000 & 1.000 & 1.000 & 1.000 \\
Persistent:
 decreasing strength & 1.000 & -- & 1.000 & 0.854 & 0.512 & 0.430 & 0.154 & 0.609 & 0.441 \\
Persistent:
 increasing strength & 1.000 & 1.000 & -- & 0.583 & 0.356 & 0.282 & 0.091 & 0.434 & 0.326 \\
Lost:
 rewiring & 1.000 & 0.854 & 0.583 & -- & 1.000 & 1.000 & 1.000 & 1.000 & 1.000 \\
Gained:
 rewiring & 1.000 & 0.512 & 0.356 & 1.000 & -- & 1.000 & 1.000 & 1.000 & 1.000 \\
Lost:
 one transient sp. & 1.000 & 0.430 & 0.282 & 1.000 & 1.000 & -- & 1.000 & 1.000 & 1.000 \\
Lost:
 both transient sp. & 1.000 & 0.154 & 0.091 & 1.000 & 1.000 & 1.000 & -- & 1.000 & 1.000 \\
Gained:
 one transient species & 1.000 & 0.609 & 0.434 & 1.000 & 1.000 & 1.000 & 1.000 & -- & 1.000 \\
Gained:
 both transient sp. & 1.000 & 0.441 & 0.326 & 1.000 & 1.000 & 1.000 & 1.000 & 1.000 & -- \\
\bottomrule
\end{tabular}
\caption{Post-hoc Dunn test (p-values). Bold values indicate significant differences (* p$<$0.05, ** p$<$0.01, *** p$<$0.001).}
\label{tab:dunn_posthoc}
\end{table}